\documentclass[../Main.tex]{subfiles}
\begin{document}

\section{Chosen Concept}
\label{chap:chosen_concept}
The previous chapters have discussed game theory and existing alternative game concepts. This knowledge resulted in the six game concepts introduced in Chapter \ref{chap:ideas}. This chapter will deliberate on the process of selecting a concept to develop. We have chosen to develop the concept \textit{Hover} introduced in section \ref{sec:hover}.

\subsection{Selection Process}
When selecting a concept, the chosen concept needs to comply with the project’s research goal, namely, developing a concept that has a socially beneficial purpose (see Section \ref{sec:research_goal}). Even though the concepts introduced in Chapter \ref{chap:ideas} are developed with social benefits in mind, some have more potential than others. We consider a concept as socially beneficial if it, to some degree, can help an individual, a group, or society as a whole. 

Secret Posts seem to have personal benefits that the player gains from moving around and being entertained by the secret posts. Likewise, Free Ticket to Ride includes individual benefits; however, it also benefits the business. Challenge Accepted and Guess Where are concepts with both individual and group benefits; however, there seem to be no benefits for society. Likewise, Puzzle with Friends focuses on group benefits, such as improved group dynamics, and collaborative skills, rather than individual benefits.  

Hover seems to be the concept that has the most substantial societal benefits. The concept motivates players to do something in the real world as individuals or as groups. It is also reasonable to argue that the concept can benefit society because it influences the local businesses or community. For instance, a player who goes to a museum benefits from gained knowledge; moreover, the museum also benefits directly or indirectly from increased income or higher attendance.

A part of the research goal is that the concept needs to be developable and testable to prove the social benefits (see Section \ref{sec:research_goal}). Additionally, due to the thesis’s time constraint, a related concern is the scope of development and testing. he introduction in Chapter \ref{chap:ideas} discussed the background for our concepts and shed light on the fact that mobile applications are highly available to users. For this reason, all the game concepts are created for the mobile platform. 

The concerns regarding development include challenges related to technical and complexity aspects. Technical concerns can be the choice of frameworks in regards to support for input and output sources. Chapter \ref{chap:chosen_tech} addresses these technical concerns for Hover, and the chosen frameworks will be presented. Complexity concerns include project duration in regards to the development scope for creating a testable prototype. Some of these concerns are disclosed under the evaluation sections in Chapter \ref{chap:ideas}. Secret Posts, for instance, may have challenges with location accuracy in multistory buildings. Puzzle with Friends may prove difficult to develop within the project duration due to the many input sensors and devices required in the real-time game. Free Ticket to Ride may also prove challenging to develop because the concept itself is based on collaborating with an existing bus company. Initiating such collaboration can become challenging if they do not want to expose their APIs. Hover will use geofencing and location tracking, which are mature technologies in most frameworks. Additionally, the concept’s scope can be scaled up if needed; however, at least a minimum viable product (MVP) seems to be within the project scope.

Although the game concepts will be equally available on mobile platforms for test subjects, other testability concerns must be recognized. Testability considerations may pose questions and concerns such as: 

\begin{itemize}
    \item What should be tested?  (User Experience/User Interface, Player Experience, behavior, enjoyment, reward mechanisms, software, etc.)
    \item What is the appropriate duration of testing?  (minutes, hours, days, weeks, etc.)
    \item How should the testing be conducted? 
    \item How will data be accumulated? (questionnaires, recorded usage data, etc.)
    \item How can data and feedback be analyzed? (quantitatively or qualitatively)
    \item How can we acquire test subjects within the target group?
    \item How do we ensure demographic variation in the test subjects?
    \item Who is the target group?
\end{itemize}

In regards to our concepts, potential testing could be performed very differently. Secret Posts, for instance, should require a multitude of users that can post and consume content. The testing should study enjoyment and network effects through techniques like interviews and questionnaires. Challenge Accepted and Guess Where are games that should be tested in a longer period, such as weeks or even a month. Data can be accumulated from questionnaires and observations and should measure effects on enjoyment, motivation, and player experience, preferably on multiple groups of friends. Free Ticket to Ride should also be tested over a greater period to see the long term effect. A questionnaire would presumably collect quantitative data and feedback that may be valuable for the business from the business’s perspective. Puzzle with Friends should be tested through observations of multiple sessions with groups of a specific size. To measure player experience, one can use techniques like questionnaires and interviews. Observing the collaborating game-play could, for instance, give insight on player enjoyment through studying facial expressions and body language. Hover should be tested over weeks, preferably on players with demographic differences. The test subjects should preferably have some preexisting relationships to be able to measure the social reward mechanisms, such as likes and comments. Data can be accumulated through questionnaires and interviews.

The evaluation of the concepts in Chapter \ref{chap:ideas} reveals that several are promising considering their testability, developability, and social benefits. The evaluation of Hover suggests that the concept seems especially interesting, and it appeals to us for several reasons. Due to the ongoing pandemic, the concept seems like a fun game to play as it encourages activities that can improve the player’s well-being. Additionally, the concept has the potential to grow in functionality and features. If the game proves successful, it could be interesting to explore and develop the concept further.


\subsection{Summary}
In this chapter, we have deliberated on selecting what game concept to develop and test. The selection process included studying the concepts’ social benefits and addressing concerns for development and testing. Out of the six, the evaluation suggests that Hover fulfills the research goal of being a socially beneficial alternative game that can be developed and tested. Lastly, we addressed personal motivations for the choice. The next chapter will go into depth about the concept and apply game theory to the concept.


\biblio % Needed for referencing to working when compiling individual subfiles - Do not remove
\end{document}